\documentclass[11pt]{article}
\usepackage[margin=1in]{geometry}
\usepackage{enumitem}
% =============================================================================
% Preambles/header.tex — shared LaTeX/Beamer preamble for this project
%
% Use from a Beamer lecture by prepending:
%     \input{header}
% and compiling with TEXINPUTS=../Preambles:$TEXINPUTS (the /compile-latex
% skill does this automatically).
%
% PALETTE CONTRACT. Color names defined here MUST match the names in
% Quarto/theme-template.scss so Beamer and Quarto renderings use the same
% palette. The scripts/check-palette-sync.sh script enforces this.
%
% To customize: edit the HEX values below AND the matching $variable in
% Quarto/theme-template.scss, then run scripts/check-palette-sync.sh.
% =============================================================================

% -----------------------------------------------------------------------------
% Engine detection + encoding
%
% Our /compile-latex skill uses XeLaTeX, where inputenc is unnecessary and
% can produce encoding warnings. Only load inputenc under pdfTeX.
% -----------------------------------------------------------------------------
\usepackage{iftex}
\ifPDFTeX
  \usepackage[utf8]{inputenc}
\fi

\usepackage{amsmath, amssymb, amsthm}
\usepackage{booktabs}
\usepackage{graphicx}

% -----------------------------------------------------------------------------
% PALETTE — mirrors Quarto/theme-template.scss variable names.
% Load xcolor and define colors BEFORE anything that references them
% (notably hyperref, hypersetup, and the Beamer color assignments).
% -----------------------------------------------------------------------------
\usepackage{xcolor}

\definecolor{primary-blue}{HTML}{012169}      % headings, accents
\definecolor{primary-gold}{HTML}{B9975B}      % emphasis, borders
\definecolor{highlight-yellow}{HTML}{F2A900}  % markers, alerts
\definecolor{light-bg}{HTML}{E8EDF5}          % title / section backgrounds
\definecolor{jet}{HTML}{1A1A1A}               % body text

% Semantic colors (used in diagrams and annotations)
\definecolor{positive}{HTML}{15803D}          % good / observed / identified
\definecolor{negative}{HTML}{B91C1C}          % bad / problematic / bias
\definecolor{neutral}{HTML}{525252}           % reference / context

% Highlight accents (match .hi-* classes in the SCSS)
\definecolor{hi-slate}{HTML}{314F4F}
\definecolor{hi-green}{HTML}{2E8B57}
\definecolor{hi-red}{HTML}{C41E3A}

% Semantic aliases so skill templates and snippets can write
% \textcolor{accent}{...} without hard-coding "primary-blue".
\colorlet{accent}{primary-blue}
\colorlet{accent2}{primary-gold}

% -----------------------------------------------------------------------------
% Hyperref — loaded AFTER xcolor + palette so link colors resolve.
% -----------------------------------------------------------------------------
\usepackage{hyperref}
\hypersetup{colorlinks=true, linkcolor=primary-blue, urlcolor=primary-blue,
            citecolor=primary-blue, pdfborder={0 0 0}}

% -----------------------------------------------------------------------------
% Course code — set once per deck (after \input{header}, before \begin{document})
% via \coursecode{CS401}. Shown in the footer on every slide so a deck's course
% is identifiable without opening the title page. Empty by default.
% -----------------------------------------------------------------------------
\makeatletter
\newcommand{\coursecode}[1]{\gdef\@coursecodetext{#1}}
\gdef\@coursecodetext{}
\makeatother

% -----------------------------------------------------------------------------
% Beamer theme assignments (only apply under Beamer).
%
% We detect Beamer via \@ifclassloaded{beamer}, which is the documented,
% non-internal way. This must be wrapped in \makeatletter/\makeatother
% because \@ifclassloaded contains an @-catcode character.
% -----------------------------------------------------------------------------
\makeatletter
\@ifclassloaded{beamer}{%
  \usefonttheme{professionalfonts}
  \setbeamercolor{structure}{fg=primary-blue}
  \setbeamercolor{frametitle}{fg=primary-blue}
  \setbeamercolor{title}{fg=primary-blue}
  \setbeamercolor{subtitle}{fg=primary-gold}
  \setbeamercolor{author}{fg=primary-blue}
  \setbeamercolor{institute}{fg=neutral}
  \setbeamercolor{date}{fg=primary-gold}
  \setbeamercolor{itemize item}{fg=primary-blue}
  \setbeamercolor{itemize subitem}{fg=primary-gold}
  \setbeamercolor{alerted text}{fg=primary-gold}
  \setbeamercolor{block title}{fg=white, bg=primary-blue}
  \setbeamercolor{block body}{bg=light-bg}
  % Semantic block variants: block=definition/concept (blue), exampleblock=
  % worked example (green/positive), alertblock=key takeaway or caution (gold).
  % Keep to <=2 colored boxes per slide across ALL three variants (INV-7).
  \setbeamercolor{block title example}{fg=white, bg=positive}
  \setbeamercolor{block body example}{bg=light-bg}
  \setbeamercolor{block title alerted}{fg=white, bg=primary-gold}
  \setbeamercolor{block body alerted}{bg=light-bg}

  % Minimal footer: course code (left) + page X / N (right), no navigation clutter
  \setbeamertemplate{navigation symbols}{}
  \setbeamertemplate{footline}{%
    \color{neutral}\footnotesize\hspace{0.7em}\@coursecodetext\hfill
    \insertframenumber/\inserttotalframenumber\hspace{0.7em}\vspace{0.3em}}
}{}
\makeatother

% -----------------------------------------------------------------------------
% TikZ setup — libraries the snippet gallery uses
% -----------------------------------------------------------------------------
\usepackage{tikz}
\usetikzlibrary{arrows.meta, positioning, calc, decorations.pathreplacing,
                fit, shapes.geometric, backgrounds}

% Shared TikZ styles — snippets and hand-written diagrams can pick these up
% via `\begin{tikzpicture}[every node/.style={...}]` or by naming specific
% styles (e.g., `\node[dag-node] (X) at (0,0) {X};`).
\tikzset{
  % Node styles for causal DAGs and similar
  dag-node/.style = {
    draw=primary-blue, fill=light-bg, rounded corners=2pt,
    minimum width=1.2cm, minimum height=0.9cm, align=center, font=\small
  },
  decision-node/.style = {
    draw=primary-gold, fill=white, diamond, aspect=1.8,
    minimum width=1.8cm, minimum height=1.2cm, align=center, font=\small,
    inner sep=1pt
  },
  % Edge styles — observed vs counterfactual
  observed-edge/.style = {-{Latex[length=2.5mm]}, thick, draw=jet},
  counterfactual-edge/.style = {-{Latex[length=2.5mm]}, thick, draw=neutral, dashed},
  confound-edge/.style = {-{Latex[length=2.5mm]}, thick, draw=negative, dashed},
  % Data-point styles for plots
  observed-dot/.style = {fill=primary-blue, draw=primary-blue, circle, inner sep=1.2pt},
  counterfactual-dot/.style = {fill=white, draw=neutral, circle, inner sep=1.2pt}
}

% -----------------------------------------------------------------------------
% Convenience macros
% -----------------------------------------------------------------------------
\newcommand{\muted}[1]{\textcolor{neutral}{#1}}
\newcommand{\key}[1]{\textcolor{primary-gold}{\textbf{#1}}}
\newcommand{\good}[1]{\textcolor{positive}{#1}}
\newcommand{\bad}[1]{\textcolor{negative}{#1}}

% Standout frame macro: full-bleed dark background for section transitions.
% Beamer-only (uses \begin{frame}/\setbeamercolor/\usebeamerfont) -- do not
% call from an article-class file (e.g. Notes/<CODE>/*-notes.tex). \muted,
% \key, \good, \bad, and \coursecode above are plain xcolor/gdef and are
% safe to use from either class; verified when Notes/article-class support
% was added.
% Expand: \transitionslide{Your transition title}
\newcommand{\transitionslide}[1]{%
  \begingroup
  \setbeamercolor{background canvas}{bg=primary-blue}
  \begin{frame}[plain]
    \centering
    \vfill
    {\usebeamerfont{title}\color{white}\Large #1\par}
    \vfill
  \end{frame}
  \endgroup
}

% Section divider frame (Beamer-only), an alternative to \transitionslide with a
% darker two-line style: near-black (jet) background, a small white label line
% above a larger gold title, and a thin gold rule along the bottom edge.
% Expand: \sectiondivider{Section 2}{Pointers}
\newcommand{\sectiondivider}[2]{%
  \begingroup
  \setbeamercolor{background canvas}{bg=jet}
  \begin{frame}[plain]
    \vspace*{3.0cm}
    {\color{white}\ttfamily\large #1\par}
    \vspace{0.5cm}
    {\color{primary-gold}\LARGE #2\par}
    \begin{tikzpicture}[remember picture, overlay]
      \draw[primary-gold, line width=2.5pt]
        ($(current page.south west)+(0,0.1)$) -- ($(current page.south east)+(0,0.1)$);
    \end{tikzpicture}
  \end{frame}
  \endgroup
}

% -----------------------------------------------------------------------------
% Channel watermark — "Concepts That Click" (YouTube).
%
% Article-class files (CompetitiveExam Q/A sets, Notes, Assignments) get a
% light diagonal draft watermark on every page plus a centered footer naming
% the channel. Beamer decks get a subtle TikZ background watermark instead
% (the footline already carries the course code). Centralized here so every
% MCQ PDF is branded without per-file edits.
% -----------------------------------------------------------------------------
\makeatletter
\@ifclassloaded{article}{%
  \usepackage{draftwatermark}
  \SetWatermarkText{Concepts That Click}
  \SetWatermarkScale{1}
  \SetWatermarkFontSize{2cm}
  \SetWatermarkLightness{0.86}
  \SetWatermarkAngle{45}
  \usepackage{fancyhdr}
  \pagestyle{fancy}
  \fancyhf{}
  \fancyfoot[C]{\footnotesize\textcolor{neutral}{Concepts That Click~|~YouTube: \href{https://www.youtube.com/@concepts.thatclick}{@concepts.thatclick}}}
  \renewcommand{\headrulewidth}{0pt}
  \renewcommand{\footrulewidth}{0pt}
  \fancypagestyle{plain}{%
    \fancyhf{}%
    \fancyfoot[C]{\footnotesize\textcolor{neutral}{Concepts That Click~|~YouTube: \href{https://www.youtube.com/@concepts.thatclick}{@concepts.thatclick}}}%
    \renewcommand{\headrulewidth}{0pt}%
    \renewcommand{\footrulewidth}{0pt}%
  }
}{}
\@ifclassloaded{beamer}{%
  \setbeamertemplate{background}{%
    \begin{tikzpicture}[remember picture, overlay]
      \node[rotate=45, opacity=0.055, align=center]
        at (current page.center)
        {\Huge\textcolor{neutral}{Concepts That Click}};
    \end{tikzpicture}%
  }
}{}
\makeatother 
\coursecode{JKSSB}
\renewcommand{\thesection}{52.\arabic{section}}
\newtheorem{example}{Example}[section]
\newenvironment{definitionbox}[1]{%
  \par\noindent\textbf{Definition (#1).}\ \itshape
}{\par\vspace{0.3em}}
\title{OS PYQ Traps and System Utilities --- Detailed Lecture Notes}
\author{JKSSB Computer Awareness}
\date{\today}

\begin{document}
\maketitle

\section{Why this lesson is a trap drill}
The computer section of the JKSSB papers does not only test whether a
candidate knows a term. It tests whether the candidate can tell a term apart
from its \emph{nearest neighbour}. A typical wrong option is not a random
word; it is the close confusion that ``sounds right'' until the exact wording
is checked. Device driver versus hardware, DMA versus ordinary CPU copying,
system software versus application software, and the Windows Run command
\texttt{winword} are all examples of this pattern.

Because of this, the present lesson is a \textbf{trap drill}. Every question
in the practice set is built around a real near-neighbour or defective-option
pattern catalogued in the course knowledge base. The aim is not to learn a
large amount of new material, but to build the habit of reading each option
as a \emph{claim to be tested}. The rest of these notes explain the small set
of facts the traps are built on, one distinction at a time.

\begin{definitionbox}{Trap drill}
A study method in which each practice option is treated as a factual claim,
and the candidate decides whether the claim is exactly correct, closely
correct but wrong, or plainly wrong.
\end{definitionbox}

\section{The trap catalogue for this lesson}
The traps this lesson drills on are collected below. Each row states the
tempting wrong belief, the correct distinction that defeats it, and the
knowledge-base reference where the pattern is recorded. Read the table as a
checklist: if an option matches a row on the left, it is the trap, and the
right-hand column tells you what to say instead.

\begin{center}
\begin{tabular}{p{0.32\textwidth}p{0.42\textwidth}p{0.10\textwidth}}
\toprule
\textbf{Tempting wrong belief} & \textbf{Correct distinction} & \textbf{Ref.} \\
\midrule
A device driver is hardware. & A device driver is software (system software); the device is hardware. & TRAP-10 \\
DMA makes the CPU copy data ``faster''. & DMA moves data between an I/O device and main memory, bypassing the CPU. & \S 52.5 \\
The OS runs from secondary storage. & While running, the OS resides in primary memory (RAM). & PYQ-22 \\
Utility software is application software. & Utility software maintains the system; application software does the user's task. & \S 52.3 \\
MS Word is system software. & MS Word is application software. & \S 52.3 \\
MS Access is an operating system. & MS Access is a database application, not an OS. & \S 52.7 \\
Windows or macOS is a Linux distribution. & Ubuntu, Fedora, and Debian are Linux distributions. & \S 52.7 \\
\texttt{winword} opens WordPad or Notepad. & \texttt{winword} opens MS Word. & \S 52.9 \\
A component is internal or external by importance. & The classification follows position relative to the system unit. & \S 52.8 \\
An RTOS is chosen for office or graphics work. & An RTOS is chosen for fixed time deadlines. & \S 52.6 \\
\bottomrule
\end{tabular}
\end{center}

\section{The operating system: role and classification}
The \textbf{operating system (OS)} is the master \textbf{system software} of a
computer. It sits between the user (and the application programs) on one side
and the hardware on the other. Its two central jobs are to provide a
\emph{user-friendly interface} and to manage the machine's resources.

The interface may be a graphical one, such as the Windows desktop, or a
command line, such as the MS-DOS prompt or the Linux shell. Either way, the
user does not need to drive the hardware directly. In its resource-management
role, the OS schedules the processor, allocates main memory, organises
storage, and controls input and output devices from a single set of
instructions. It also provides services --- opening a file, printing a
document, sending data over a network --- that application programs call
instead of talking to the hardware themselves.

\begin{definitionbox}{Operating system}
The system software that controls the computer's hardware and provides a
user-friendly interface and a set of services to application programs.
\end{definitionbox}

Common operating systems are Windows, Linux, UNIX, MS-DOS, Android, and
macOS. An application program is not an operating system. MS Word, MS Excel,
MS Access, a web browser, and a media player are applications, not operating
systems. This is one of the exact-recognition distinctions the exam uses.

A frequently tested additional fact is \emph{where} the OS is when the
computer is running. The OS is stored permanently on secondary storage (the
disk), but while the machine is running it is loaded into \textbf{primary
memory (RAM)}. The CPU executes instructions only from main memory, so the
running OS must be resident there. When power is removed, the RAM copy is
lost and the OS is loaded again from disk at the next boot.

\section{System, application, and utility software}
Software is grouped by its purpose. \textbf{System software} runs and manages
the machine: the operating system, device drivers, utility programs, and
language translators (compiler, interpreter, assembler) all belong here.
\textbf{Application software} performs a task the user wants done: typing a
document, preparing a spreadsheet, browsing the web. \textbf{Utility
software} is a small maintenance category: tools such as Disk Cleanup, Disk
Defragmenter, and antivirus programs that keep the system healthy and
efficient. Utility software is closely tied to system software, but it is
useful to name it separately because exam options mix the three categories.

\begin{center}
\begin{tabular}{p{0.24\textwidth}p{0.34\textwidth}p{0.32\textwidth}}
\toprule
\textbf{Category} & \textbf{Role} & \textbf{Examples} \\
\midrule
System software & Runs and manages the machine & OS, device drivers, translators \\
Application software & Does a user's task & MS Word, Excel, a browser \\
Utility software & Maintains and optimises the system & Disk Cleanup, Defragmenter, antivirus \\
\bottomrule
\end{tabular}
\end{center}

\begin{example}
An item asks, ``Which of the following is application software?'' with the
options Windows, MS Word, Disk Defragmenter, and a device driver. Windows is
an operating system (system software); Disk Defragmenter is utility software;
a device driver is system software too. Only \textbf{MS Word} is application
software. The trap is that all four options are well-known pieces of
software, and only the \emph{category} separates them.
\end{example}

\section{Device drivers are software}
A \textbf{device driver} is a small software program that allows the
operating system to communicate with a specific hardware device. Each device
--- a printer, a keyboard, a graphics card, a network adapter --- needs a
matching driver. The driver receives general commands from the OS and
translates them into the device-specific instructions that the hardware
understands. Without the correct driver, the OS cannot use the device
properly.

The trap here is a classification trap. A device driver \emph{controls}
hardware, but the driver itself is \textbf{software}, and because it extends
the operating system's control of the machine it is counted as \textbf{system
software}. The device it drives is hardware. The OS itself is system
software. An option that says ``device driver is hardware'' has confused the
software with the device.

\begin{definitionbox}{Device driver}
A software program that lets the operating system control a particular
hardware device by translating OS commands into device-specific
instructions.
\end{definitionbox}

\section{DMA: transferring data without the CPU}
Normally, moving data between an input/output device and main memory keeps
the CPU busy: the CPU reads each unit from the device and writes it to memory
(or the reverse), one step at a time. \textbf{DMA} --- \textbf{Direct Memory
Access} --- replaces this with a dedicated transfer. A DMA controller moves
data \emph{directly} between the I/O device and main memory while the CPU is
bypassed for the duration of the transfer. The CPU starts the transfer and is
notified by an interrupt when it is complete.

\begin{definitionbox}{Direct Memory Access (DMA)}
A method of data transfer in which a controller moves data between an
input/output device and main memory directly, without the CPU handling each
individual byte.
\end{definitionbox}

The advantages are clear. Because the CPU is not occupied with copying each
byte, it is free to run other instructions, and large block transfers finish
faster. Several traps grow from this description. DMA moves data between a
\emph{device and main memory}, not between the CPU and the device. The CPU is
\emph{bypassed} --- it is not made ``faster''. DMA does not increase the size
of RAM and does not encrypt data. The CPU still initiates the transfer and
handles the completion interrupt; it is simply not in the per-byte path.

\section{Real-time operating systems}
A \textbf{real-time operating system (RTOS)} is an operating system designed
so that tasks complete within \textbf{fixed time deadlines}. In an ordinary
desktop OS, a task may take as long as the scheduler permits as long as it
finishes eventually. In an RTOS, meeting the deadline is the whole point: a
late response is a failure, even if the answer is correct.

An RTOS is therefore chosen for embedded and control systems, such as
industrial machinery, medical devices, and vehicle electronics, where a
delayed action can be dangerous. It is not chosen for typing a document,
designing a website, or editing a photograph; those are ordinary
application tasks that do not carry hard deadlines.

\begin{definitionbox}{Real-Time Operating System (RTOS)}
An operating system that guarantees that critical tasks complete within a
bounded, fixed time so that deadlines are always met.
\end{definitionbox}

\section{DOS, UNIX, and Linux distributions}
\textbf{DOS} stands for \textbf{Disk Operating System}. MS-DOS was a
command-line operating system: the user typed commands rather than clicking
icons. \textbf{UNIX} is a multi-user, multitasking operating system that
influenced many later systems. \textbf{Linux} is an open-source, UNIX-like
operating system; because its source can be freely used and modified, it is
distributed in many packaged forms called \textbf{distributions}. Well-known
Linux distributions include \textbf{Ubuntu}, \textbf{Fedora}, \textbf{Debian},
and \textbf{Linux Mint}.

This gives another exact-recognition drill. Windows and macOS are operating
systems, but they are \emph{not} Linux distributions. MS-DOS is an operating
system but, again, not a Linux distribution. Only Ubuntu, Fedora, Debian, and
Mint belong on the Linux-distribution list. Similarly, MS Access is a
database application and is \textbf{not} an operating system at all, even
though it is a familiar Microsoft product.

\section{Internal and external hardware}
Hardware is classified by its \emph{position} relative to the system unit
(the cabinet that houses the main machine). \textbf{Internal} hardware sits
inside the system unit: the motherboard, the CPU, RAM, the hard disk, the
power supply, and expansion cards. \textbf{External} hardware, also called a
\textbf{peripheral}, sits outside the system unit and is joined to it through
ports and cables: the monitor, keyboard, mouse, printer, and scanner.

The trap is that a familiar word can appear in the wrong column. A monitor
and a keyboard are external; a motherboard and a RAM module are internal. The
classification follows the position, not the importance of the component.

\begin{definitionbox}{Internal vs external hardware}
Internal hardware is located inside the system unit; external hardware
(a peripheral) is located outside it and connected through ports.
\end{definitionbox}

\section{The Windows Run box and the \texttt{winword} command}
The Windows \textbf{Run} box, opened with the Windows + R shortcut, starts a
program when the user types that program's command. The command
\texttt{winword} opens \textbf{MS Word}. In the same way,
\texttt{notepad} opens Notepad, \texttt{wordpad} opens WordPad, and
\texttt{cmd} opens the Command Prompt. The command must match the program
exactly.

\begin{example}
An item asks which program is opened by typing \texttt{winword} in the Run
box, offering Notepad, MS Word, WordPad, and File Explorer. The answer is
\textbf{MS Word}. Notepad is opened by \texttt{notepad}, WordPad by
\texttt{wordpad}, and File Explorer by \texttt{explorer}. The near-neighbours
are deliberately word-processor or editor names, so the candidate who only
half-remembers the command can be pulled to the wrong option.
\end{example}

\section{Running example: the office desktop}
It helps to hang these facts on one concrete machine. Picture the office
desktop used throughout this course: a multi-core PC running Windows, with
MS Word and MS Excel installed, a printer, and an internet connection.

When the machine is switched on, the Windows OS is loaded from the disk into
primary memory and stays resident there while the machine runs. MS Word and
MS Excel are \emph{application} software that the user opens to do work.
Disk Cleanup and Disk Defragmenter are \emph{utility} software that the user
runs occasionally to keep the machine healthy. The printer does not work
until its \emph{device driver} --- a piece of system software --- is
installed, so that Windows can translate its print commands into the
printer's own instructions. When a large print job or file copy is under way,
a \emph{DMA} transfer can move the data between the device and main memory
while the CPU continues with other work.

\begin{example}
A single scenario can therefore generate several trap options about one
machine. ``Word is system software'' is false --- Word is application
software. ``The printer driver is hardware'' is false --- the driver is
software. ``The OS is stored on the disk and run from there'' is half-true
and therefore defective: it is stored on the disk but \emph{run} from main
memory. Recognising which category each item belongs to settles all three.
\end{example}

\section{Defective options and how to read them}
Some official options are not merely wrong; they are \emph{defective} --- the
wording itself is inconsistent or ambiguous. The knowledge base records this
as a real feature of the papers (Anti-Pattern AP-02). A defective option is
teaching material, not a fact to memorise. The safe method is to quote the
option exactly, mark the defect, and then apply the stable concept.

\begin{example}
A defective option might read, ``The device driver is the hardware that
stores the instructions the operating system uses to control a device.'' The
defect is the phrase ``the hardware that stores the instructions''. A driver
is \emph{software} that holds those instructions; it is not hardware. The
stable concept to apply is TRAP-10: a device driver is software. Once the
defect is spotted, the option can be rejected without any doubt.
\end{example}

The same discipline applies to near-neighbour options that are not defective
but simply close. When two options are both plausible, compare them word by
word against the definitions in this lesson. The distinction is usually a
single term --- \emph{software} versus \emph{hardware}, \emph{device to
memory} versus \emph{CPU to device}, \emph{application} versus
\emph{utility}, or \texttt{winword} versus \texttt{wordpad}. Reading for that
one term is what turns a guess into a correct answer under negative marking.

\section{Terminology ladder for this lesson}
Direct-recall expansion items are the single most common JKSSB format, so the
exact terms from this lesson are collected here with their full forms. Each
row is one fact.

\begin{center}
\begin{tabular}{p{0.13\textwidth}p{0.30\textwidth}p{0.47\textwidth}}
\toprule
\textbf{Term} & \textbf{Full form} & \textbf{One-line meaning} \\
\midrule
OS & Operating System & System software that provides an interface and manages resources \\
DMA & Direct Memory Access & Data moves between an I/O device and main memory with the CPU bypassed \\
RTOS & Real-Time Operating System & An OS that meets fixed time deadlines \\
DOS & Disk Operating System & A command-line operating system (MS-DOS) \\
Driver & Device driver & Software that lets the OS control a specific device \\
Utility & Utility software & Maintenance software such as Disk Cleanup or Defragmenter \\
Peripheral & External hardware & A device outside the system unit, joined by ports \\
\bottomrule
\end{tabular}
\end{center}

\section{Exam retrieval and common confusions}
Six distinctions prevent most errors on this topic.
\begin{enumerate}
  \item A device driver is \textbf{software} (system software); the device it
    controls is hardware (TRAP-10).
  \item The OS is \textbf{system software}. MS Word is application software;
    Disk Cleanup is utility software.
  \item While running, the OS is held in \textbf{primary memory (RAM)}, not
    executed from secondary storage (PYQ-22).
  \item DMA moves data between a \textbf{device and main memory} and
    \textbf{bypasses the CPU}; it is not ``faster CPU copying''.
  \item \textbf{RTOS} = Real-Time Operating System, chosen for fixed
    deadlines.
  \item \textbf{DOS} = Disk Operating System; Linux distributions include
    Ubuntu, Fedora, and Debian. \texttt{winword} opens MS Word.
\end{enumerate}

\subsection*{Exercises}
\begin{enumerate}[label=\arabic*.]
  \item Classify each as system, application, or utility software: Windows,
    MS Word, Disk Defragmenter, a device driver, Disk Cleanup.
  \item Explain why a device driver is described as software even though it
    controls hardware.
  \item Describe a DMA transfer: name the two locations between which data
    moves and state the role of the CPU.
  \item State the full form of RTOS and give one reason an RTOS is used.
  \item Which of these are operating systems, and which are Linux
    distributions: Ubuntu, Windows, Debian, MS Access, Fedora?
  \item What does typing \texttt{winword} in the Windows Run box do, and how
    does it differ from \texttt{wordpad}?
\end{enumerate}

\subsection*{Solutions}
\begingroup\small
\begin{enumerate}[label=\arabic*.]
  \item System software: Windows and the device driver. Application
    software: MS Word. Utility software: Disk Defragmenter and Disk Cleanup.
  \item The driver is a stored set of instructions that the OS loads and
    runs to talk to the device. It has no physical body of its own; the
    hardware is the device. A program that extends OS control of hardware is
    therefore software, and specifically system software.
  \item Data moves directly between the input/output device and main memory.
    A DMA controller performs the transfer; the CPU starts it and is
    interrupted when it finishes, but it does not handle each byte.
  \item RTOS = Real-Time Operating System. It is used where tasks must
    complete within fixed time deadlines, as in embedded and control
    systems.
  \item Operating systems: Ubuntu (also a Linux distribution), Windows,
    Debian (also a Linux distribution), Fedora (also a Linux distribution).
    MS Access is not an operating system; it is a database application.
  \item Typing \texttt{winword} opens MS Word. Typing \texttt{wordpad} opens
    WordPad, a different text editor. The two commands open different
    programs and must not be confused.
\end{enumerate}
\endgroup
\end{document}
